\section{Goals of Usability Testing}
\label{sec:goals}

These goals shape each session and the success criteria in Section~\ref{sec:task_scenarios}. They map to the metrics in Section~\ref{sec:metrics}.

\begin{enumerate}
  \item \textbf{Onboarding and verification:} See whether participants can register with a McMaster-affiliated account, complete verification, and choose a role, \textbf{including} understanding and fixing sign-up errors (wrong email format, weak password, missing driver license when required).
  \item \textbf{Rider journey:} See whether they can find or request a ride, open trip details when needed, and \textbf{read the estimated cost} shown on match proposals.
  \item \textbf{Driver journey:} See whether they can post a ride with direction and schedule, and work with passenger requests as the app presents them.
  \item \textbf{Matching and trip state:} See whether they understand proposals, \textbf{accept or decline} matches or requests, handle ``no match'' or blocked states, and move between related screens without getting lost.
  \item \textbf{After the ride:} See whether they can \textbf{submit a review} and \textbf{submit a report} through the in-app safety flow when the moderator provides a suitable trip state.
  \item \textbf{Profile:} See whether they can \textbf{update profile information} (and optionally find settings) without undue friction.
  \item \textbf{Error recovery:} Note hesitation, wrong taps, and dead ends, and whether participants recover with at most reasonable moderator help.
  \item \textbf{Satisfaction:} Capture perceived ease and confidence with SUS and short debrief questions.
\end{enumerate}
